\section{A projected row likelihood and the spin spectral measure}
\label{sec:row-spin}

Projection onto the terminal configuration improves the moment bounds
for a row likelihood from exponential to subexponential in the path
length.  We prove this improvement first, keeping separate the roles of
the two estimates: unprojected moments justify passage to a fixed-time
limit, and projected moments determine the hidden-sector decay rate.
We then construct the spin spectral measure and derive the two strict
renewal inequalities needed to sum the hidden path expansion.

The inputs are the rough gradient and likelihood estimates of
\Cref{prop:path-rough}, the cavity density estimate of
\Cref{prop:stat-density}, and the following specified parts of
\Cref{thm:path-calculus}: deterministic stationary spin and path limits,
conditional prediction on causal trees, the noncrossing trace rule, and
the two-independent-equilibrium-sample comparison for spin histories,
fields, and likelihood labels.  The last comparison does not assert
independence of arbitrary recipes with noncentered queries or shared
random seeds.  The bridge argument below concerns an independent
\emph{removed Gaussian row}; it is not an application of a path rule to a
new bridge conditioned on two prescribed endpoints.

\subsection{The statement and its quantifiers}

Fix a row $i$ and $d>0$.  A mask $M\subset[0,d]$ is a deterministic
finite union of intervals specifying when the row perturbation is
present.  In the diagonal likelihood representation of
\Cref{prop:path-gradient}, let $\ell_i^M$ change the field at every site
$j\ne i$ by $-2X_i(s)J_{ij}$ when $s\in M$, leaving the site-$i$ update
rule unchanged.  The law with this modification has path density
$\ell_i^M$ with respect to the ordinary path law.  Put
\begin{equation}
 A_{i,x}^M(y)=\E_x[\ell_i^M\mid X_d=y].
 \label{eq:row-projection}
\end{equation}
This conditional expectation is the ratio of the modified and ordinary
terminal transition probabilities.  It includes paths on which the last
site-$i$ refresh left the spin unchanged, so the attempt and flip
likelihood conventions agree in \Cref{eq:row-projection}.

\begin{theorem}[Projected moments of a deterministic row likelihood]
\label{thm:row-projected-moments}
For each fixed real $r$ and each $\varepsilon>0$, there is
$d_0=d_0(r,\varepsilon)$ such that, for every fixed $d\ge d_0$ and every
deterministic mask $M$, outside a disorder event of probability tending
to zero there is a set $B_{n,d,M}$ satisfying
\begin{equation}
\mu(B_{n,d,M})\le e^{-c n},\qquad
 \max_i\E_x\bigl[(A_{i,x}^M(X_d))^r\bigr]
       \le e^{\varepsilon d}+o_n(1)
 \quad(x\notin B_{n,d,M}).
 \label{eq:row-moment-quantified}
\end{equation}
Here $c>0$ can depend on the fixed $d,r,\varepsilon$, but the estimates
in the proof are uniform in the deterministic mask.  Thus the notation
$e^{o_r(d)}$ for these moments always means dimension first and then
$d\to\infty$.  For $0\le r\le1$ the upper bound $1$ follows directly
from Jensen's inequality.  Negative $r$ in
\Cref{eq:row-moment-quantified} refers to the unrestricted positive
likelihood, not to a likelihood multiplied by a cutoff indicator.

The theorem also holds in its integrated form: if $M=M_\alpha$ ranges
over a deterministic measurable family with a finite measure on its
parameter space, the good disorder event can be chosen so that the
\emph{integrated} Gibbs mass of the exceptional starts is exponentially
small.  This includes masks determined by simplex times in a fixed
finite path expansion and the independent Poisson schedules used below.
It does not assert a simultaneous bound for a mask selected to maximize
an error after seeing the disorder.
\end{theorem}

The maximum over $i$ is obtained by a union bound over the $n$ rows,
applied after the exponential estimate for a row and a start.  This
exponential margin also allows any fixed polynomial weight and any
initial density bounded above by $e^{o(n)}$ in the integrated assertion.
We will use precisely these quantifiers in \Cref{thm:hidden-propagation}.

\subsection{The cavity identity and the last refresh}

We begin by expressing the projected likelihood through a cavity chain.
Remove row and column $i$, and let $Y$ be the resulting chain started
from $x_{-i}$.  Conditional on the cavity disorder, the removed vector
$(J_{ij})_{j\ne i}$ is independent Gaussian with covariance
$\beta^2 I/n$.  Let $N_j$ count the rate-one refresh attempts at $j$ and
let
\[
 M_j(t)=\sum_{s\le t:\,\text{refresh at }j}
       \bigl(Y_j(s)-m_j^Y(s-)\bigr),\qquad
 m_j^Y=\tanh((J^{(i)}Y)_j).
\]
These martingales are orthogonal, with predictable quadratic variations
$\int_0^t v_j^Y(s)\dd s$.  In the likelihood expansion we also retain
the clock-conditioned compensator $\int_0^t v_j^Y(s-)\dd N_j(s)$.

Let $\mathcal S$ be the site-$i$ refresh schedule. First condition on
a deterministic realization $\mathfrak s=(t_1,\ldots,t_m)$, with
$0<t_1<\cdots<t_m<d$.  The added spin $a$ stays constant between
successive scheduled times.  At a refresh its new value is sampled with
probability
\[
 p_\sigma(H)=\frac{e^{\sigma H}}{2\cosh H},
 \qquad H=\sum_{j\ne i}J_{ij}Y_j,
 \qquad\sigma\in\{-1,1\}.
\]
Set $\theta=1-2\1_M$.  Relative to the cavity experiment with these
site-$i$ coins, adding the field $\theta(s)a(s)J_{ij}$ at the other sites
has likelihood $L_\theta$.  Taylor expansion of $\log\cosh$ gives
\begin{align}
 \log L_\theta
 &=\sum_{h=0}^{m}\left(a_h Z_h-\tfrac12 B_h\right)+o_{\P}(1),
 \label{eq:row-score-expansion}\\
 Z_h&=\sum_{j\ne i}J_{ij}
                \int_{I_h}\theta(s)\dd M_j(s),\qquad
 B_h=\sum_{j\ne i}J_{ij}^2
                \int_{I_h}v_j^Y(s-)\dd N_j(s),
 \nonumber
\end{align}
where $I_h$ are the intervals between scheduled refreshes.  All changes
of the mask within $I_h$ remain inside a single stochastic integral.
Allowing arbitrarily many such changes therefore does not enlarge the
Gaussian list.  The absolute Taylor remainder is at most
\[
 C\max_{j\ne i}|J_{ij}|
       \sum_{j\ne i}J_{ij}^2N_j(d).
\]
At fixed $d$, this bound tends to zero in probability and admits the
moment truncations used below.  The time remains fixed before the
dimension limit throughout this expansion; it is never chosen as a
function of the dimension.

To keep track of the terminal spin, force its last refresh to take the
value $\sigma$, leaving the probabilities of all earlier coins
unchanged.  With $s_*=t_m$, define
\[
 \bar H_* =J_i\E_{x_{-i}}Y(s_*),\qquad
 q_\sigma=p_\sigma(\bar H_*),\qquad
 F_\sigma=\frac{p_\sigma(H(s_*))}{p_\sigma(\bar H_*)}.
\]
The elementary inequality $|\partial_H\log p_\sigma(H)|\le2$ gives
\begin{equation}
 e^{-2|H(s_*)-\bar H_*|}\le F_\sigma
                  \le e^{2|H(s_*)-\bar H_*|}.
 \label{eq:row-forced-coin}
\end{equation}
This ratio belongs inside the cavity average: it restores the random
last-spin probability after the deterministic reference probability has
been factored out.  Set $Z=Y(d)$ and define
\[
 A_{\theta,\sigma}(Z)
   =\E_{\mathrm{cav}}[L_\theta F_\sigma\mid Z],
\]
where the expectation includes the earlier site-$i$ coins and the
forced last coin.  When no refresh occurs, only one final spin is admissible, and
we set $q=1$ and $F=1$.  If $Q_d$ is the cavity terminal law, this
construction gives the exact kernel identities, conditional on
$\mathcal S=\mathfrak s$,
\begin{equation}
 \begin{aligned}
 P_\theta(X_i(d)=\sigma,Y(d)\in\dd z)
       =q_\sigma A_{\theta,\sigma}(z)Q_d(\dd z),
 \\
 A_{i,x}^{M,\mathfrak s}(\sigma,z)
       =\frac{A_{\theta,\sigma}(z)}{A_{1,\sigma}(z)}.
 \end{aligned}
 \label{eq:row-cavity-ratio}
\end{equation}
Here $A_{i,x}^{M,\mathfrak s}(y)=
\E_x[\ell_i^M\mid X_d=y,\mathcal S=\mathfrak s]$; the schedule
is still observed. The $r$th moment under the ordinary terminal law
conditional on that schedule is therefore
\begin{equation}
 \sum_\sigma q_\sigma
  \E_{Q_d}\!\left[A_{\theta,\sigma}^{\,r}
                         A_{1,\sigma}^{\,1-r}\right].
 \label{eq:row-ratio-moment}
\end{equation}
We have reduced the desired bound to fixed positive and negative
moments of two conditional cavity averages.  Their weights $q_\sigma$
sum to one even when the deterministic mean field $\bar H_*$ is large.

\begin{lemma}[Fixed-time moment passage]
\label{lem:row-uniform-integrability}
Fix $d$, a schedule, a mask, and $q<\infty$.  On the disorder events of
\Cref{prop:path-rough}, the positive and negative moments of $L_\theta$
are bounded by $e^{C_qd}$, uniformly in the schedule and mask.  The
variables $L_\theta F_\sigma$ have finite moments of every fixed positive
and negative order, with constants uniform in dimension at this fixed
$d$.  The same is true of their conditional averages given $Z$.
Bounded replicated-bridge tests, followed by these bounds, determine
both the positive and negative moments of the conditional averages in
the dimension limit.
\end{lemma}

\begin{proof}
Consider first a refresh at another site.  The logarithmic perturbation
is a Bernoulli likelihood increment with field displacement
$\theta aJ_{ij}$.  Its conditional exponential moment, at any fixed
real order, is at most
$\exp(C_qJ_{ij}^2)$.  Iteration over the attempts, followed by the Poisson
exponential formula, gives $e^{C_qd}$ because
$\sum_jJ_{ij}^2\le C$ and $\max_j|J_{ij}|=o(1)$.  The same calculation controls inverse moments.  After its prescribed
refresh the forced spin is predictable, so it does not change the
calculation.

We next control the forced-coin factor.  For $f(y)=J_iy$,
\Cref{prop:path-rough} supplies a fixed-horizon bound
on $\sup_y\norm{dS_sf(y)}^2$.  Applying
\Cref{lem:input-mgf} controls every fixed exponential moment of
$f(Y(s_*))-\E f(Y(s_*))$.  Combine this with
\Cref{eq:row-forced-coin} and H\"older's inequality.  The resulting constant may grow rapidly with fixed $d$.  This step
provides uniform integrability at fixed time, rather than a large-$d$
rate.

For a positive $V$ and $A=\E[V\mid Z]$, Jensen gives
\begin{equation}
 A^q\le\E[V^q\mid Z]\quad(q\ge1),\qquad
 A^{-q}\le\E[V^{-q}\mid Z]\quad(q>0).
 \label{eq:row-jensen-ui}
\end{equation}
Choosing an order strictly above the desired one gives uniform
integrability for the required powers of $A$.  To identify those powers
by replicated tests, take copies $V^{(1)},\ldots,V^{(b)}$ that are
conditionally independent given $Z$.  They satisfy
\[
 \E\left|b^{-1}\sum_{a=1}^bV^{(a)}-A\right|^2
       \le b^{-1}\E[V^2].
\]
Cap the $V^{(a)}$ first and pass to the dimension limit with the cap
and $b$ fixed.  Then let $b\to\infty$ and remove the cap.  The higher
moments in \Cref{eq:row-jensen-ui} justify this passage for positive and
negative powers alike; in particular, a limiting $A$ has no mass at
zero.  H\"older's inequality in \Cref{eq:row-ratio-moment} transfers the
conclusion to the terminal density ratio.
\end{proof}

\subsection{Replicated Gaussian bridges}

The terminal projection couples cavity paths through a common endpoint.
We represent that coupling by a shared Gaussian part, then show that
scores accumulated before a terminal buffer have a small shared part.
Their remaining private part can then be integrated in chronological
order. To construct the representation, fix the time $d$, the schedule, the moment
orders, and all preliminary truncations before taking the dimension
limit; $d$ will grow only later.  For each terminal state $Z$, take a
fixed number of conditionally independent cavity bridges from the same
specified start to $Z$.  On each branch, form the vector list
\begin{equation}
 \begin{aligned}
 \frac{Y(s)-\E Y(s)}{\sqrt n}\quad
       &\text{at the scheduled coin and last-field times},\\
 \frac1{\sqrt n}\int_{I_h}\theta\dd M
       \quad&\text{on the required score intervals}.
 \end{aligned}
 \label{eq:row-bridge-list}
\end{equation}
Conditional on these cavity data, projection by the independent removed
row gives an exactly Gaussian list, with covariance $\beta^2$ times
its Gram matrix.  The mean fields require a separate step: replace
$J_i\E Y(s)$ in the coin probabilities by formal deterministic
parameters $\bar h_s$.  We establish the averaged-test comparisons and
bounds uniformly in those parameters, and only then substitute the
actual values $\bar h_s=J_i\E Y(s)$.  This order avoids conditioning
the row on its own mean-field projections while continuing to use its
original Gaussian covariance.

To obtain the quenched statement, we now concentrate the averaged tests
in the removed row.  The required gradient bounds come from covariance
operators under the ordinary cavity path law.  For each unscaled
centered field vector in \Cref{eq:row-bridge-list}, that operator is
bounded by a constant depending only on $d$: apply the variance
conclusion of
\Cref{lem:input-mgf} to every deterministic linear combination of spins,
using \Cref{prop:path-rough}.  For an unscaled score vector the covariance
operator is bounded by $d$, directly by orthogonal martingale isometry.
Although the branches share an endpoint, each has the ordinary path law
as its unconditional marginal.  For a bounded smooth function $G$ of a
fixed list of row projections, Cauchy--Schwarz with the covariance
bounds therefore controls the Euclidean gradient of
$\E_{\mathrm{cav}}G$ by a constant depending only on the fixed list.
Since the row covariance is $\beta^2 I/n$, Gaussian concentration gives
error probability $e^{-c n}$ at each fixed accuracy.

We apply the same procedure to the brackets, first capping each site's
count.  With the cap fixed, concentration of a quadratic Gaussian form
gives
\begin{equation}
 B_h-\frac{\beta^2}{n}\sum_{j\ne i}
                   \int_{I_h}v_j^Y\dd N_j\longrightarrow0.
 \label{eq:row-bracket-replacement}
\end{equation}
Poisson tails remove the count cap from the averaged test, and the
moment bounds of \Cref{lem:row-uniform-integrability} then permit the
required moment passage.  The Taylor remainder in
\Cref{eq:row-score-expansion} is handled in the same order.  The site-$i$ coins are finite smooth
functions of the listed fields; after capping the tested likelihoods,
the same concentration estimate applies to their averages.  For uniformity in the deterministic mean fields, use a finite net on a
compact field range and then the limits of $p_\sigma(h+u)$ as
$h\to\pm\infty$, truncating the centered $u$ first.  This last step
allows the deterministic mean field to be unbounded.

It remains to transfer these estimates from a cavity start to a full
Gibbs start.  They are uniform over the start for every good cavity
matrix, and the two Gibbs laws satisfy
\[
 \frac{\mu_J(a,y)}{\mu_{J^{(i)}}(y)}
    =\frac{e^{aJ_iy}}{2\E_{\mu_{J^{(i)}}}\cosh(J_iY)}.
\]
Since $\norm{J_i}\le C$, the logarithm of this ratio has magnitude
$O(\sqrt n)$.  The change of measure therefore preserves each fixed exponential
failure rate.  A union bound over the rows followed by Markov's
inequality gives exponentially small quenched Gibbs mass of bad starts
on an event of disorder probability tending to one.  At this stage the
Gaussian description concerns bounded tests; its extension to moments
uses \Cref{lem:row-uniform-integrability}.

This Gaussian description does not require deterministic self-Grams.
The following representation retains their randomness while separating
the part shared by different bridges from the private Gaussian part of
each bridge. This is a specialized Dovbysh--Sudakov-type construction;
for the general representation, see Panchenko~\cite[Proposition~1,
equation~(1.2), and Lemmas~2--3]{Panchenko2010DS}. The finite-list
principal-coordinate proof needed here is given below, rather than
invoking the general representation theorem as a black box.

\begin{lemma}[Shared and private Gaussian parts]
\label{lem:row-hilbert-representation}
After subsequence extraction, a replicated array of the lists in
\Cref{eq:row-bridge-list} has the following representation.  Conditional
on a directing variable $\mathcal D$, each branch draws independently
its self-Gram matrix $Q$, its bracket data, and a list of vectors
$r_\alpha$ in one separable real Hilbert space.  Cross-Grams between
distinct branches are $\ip{r_\alpha}{r'_\gamma}$, and
$Q-(\ip{r_\alpha}{r_\gamma})_{\alpha,\gamma}$ is positive
semidefinite.  Include the factor $\beta$ in the vectors.  Conditional
on these data, the row projections have the form
\begin{equation}
 g(r_\alpha)+\xi_\alpha,
 \label{eq:row-shared-private}
\end{equation}
where $g$ is one common isonormal process and, independently on each
branch, $\xi$ is centered Gaussian with covariance
$Q-(\ip{r_\alpha}{r_\gamma})$.  Conditional averages in this
representation integrate one branch and its coins with $\mathcal D,g$
held fixed.
\end{lemma}

\begin{proof}
Begin with a finite-dimensional list and condition on $Z$.
Diagonalize
\[
 T=\E\left[\sum_\alpha u_\alpha\otimes u_\alpha\mid Z\right],
\]
where $u_\alpha$ denotes a normalized list vector.  Work first on the
sets where the list norms are bounded.  If $\Pi_k$ projects on the first
$k$ eigenvectors, independence of two branches and
$\Tr T\le K$ give
\[
 \E\left[\left|\ip{(I-\Pi_k)u_\alpha}
                           {(I-\Pi_k)u'_\gamma}\right|^2\Bigm|Z\right]
       \le\Tr[((I-\Pi_k)T)^2]\le\frac{K^2}{k+1}.
\]
Extract the laws of the finite coordinate lists together with their
self-Grams, and then let $k\to\infty$.  Count tails and the spin-norm
bounds remove the preliminary restriction on list norms.  The resulting
directing law on lists in $\ell^2$ identifies the limiting cross-Grams.
At every finite coordinate projection, the stated covariance difference
is positive semidefinite, and this property persists in the limit.  A common isonormal
process on $\ell^2$, plus independent Gaussians with the covariance
difference, realizes exactly these Gaussian Gram arrays.  The
conditional-average assertion follows from the replicated
approximation in \Cref{lem:row-uniform-integrability}.
\end{proof}

\subsection{The shared part of a prefix score}

The directing-law representation retains the information shared by
bridges with the same endpoint. We now show that scores accumulated
well before the terminal time have very little shared Gaussian part.  For this \emph{prefix score}
estimate, set $s=\sqrt d$ and $T=d-s$.  Split scores at $T$ when
necessary, leaving at most $m+2$ intervals.  The intervening suffix
allows us to use the cavity density bound.  Write $\kappa(s)$ for the
constant from the cavity version of \Cref{prop:stat-density}, so that
\[
 P_s^{\mathrm{cav}}(y,z)\le e^{n\kappa(s)}\mu_{\mathrm{cav}}(z),
 \qquad \kappa(s)\le C e^{-c s^{2/3}}.
\]
For any $a>0$, discard the endpoints at which
$Q_d(z)<e^{-an}\mu_{\mathrm{cav}}(z)$.  Their $Q_d$-probability is at
most $e^{-an}$.  On the retained endpoints, the joint law of two bridge
prefixes, after the common endpoint is integrated out, is dominated by
\begin{equation}
 e^{n(2\kappa(s)+a)}
       \times\text{the law of two independent cavity prefixes}.
 \label{eq:row-bridge-domination}
\end{equation}
Indeed its density given the two prefix endpoints $y,y'$ is
$\sum_zP_s(y,z)P_s(y',z)/Q_d(z)$, restricted to the retained $z$;
insert the preceding two bounds and sum $\mu_{\mathrm{cav}}(z)$.

We first estimate overlaps under the independent-prefix law, where the
normalized inner product of any two score vectors is centered.  Stop
each coordinate martingale when its clock count reaches a fixed large
polynomial in $d$.  The resulting product martingale has jumps bounded
by a polynomial in $d$ and bracket bounded by $n$ times a polynomial in
$d$.  The bounded-jump exponential martingale inequality then gives,
for any inverse-polynomial tolerance,
\begin{equation}
 \P\left(\left|n^{-1}\sum_j
        U_j^{(1)}U_j^{(2)}\right|>d^{-b}\right)
           \le 2e^{-n/d^{c_b}}.
 \label{eq:row-independent-overlap}
\end{equation}
Constants can absorb the fixed number of score positions.  Removing
the coordinate stopping changes normalized squared norms by an
arbitrarily small inverse power of $d$: the relevant bound is a
constant times $n^{-1}\sum_j N_j(d)^2\1_{N_j(d)>d^C}$, whose expectation
has this property and whose spatial average converges in probability.
These removals are made on each \emph{single bridge marginal}, which
is the ordinary path law, before using
\Cref{eq:row-bridge-domination}; a merely subexponential count-tail
bound is never multiplied by the bridge density.
Since $\kappa(\sqrt d)$ is smaller than every inverse power of $d$,
choose $a$ with the same property and apply
\Cref{eq:row-bridge-domination} to
\Cref{eq:row-independent-overlap}.  For all sufficiently large fixed
$d$, every prefix-score cross-Gram is bounded by an arbitrarily
prescribed inverse power of $d$ almost surely in the dimension limit.

The almost-sure bound on cross-Grams is the key output of this
comparison.  The next lemma converts it into a bound on shared-vector
norms; a small mean squared overlap would not suffice.

\begin{lemma}[An almost-sure overlap bound in a separable space]
\label{lem:row-overlap-norm}
If $r,r'$ are independent with the same law in a separable Hilbert
space and $|\ip r{r'}|\le\delta$ almost surely, then
$\norm r^2\le\delta$ almost surely.  The same assertion holds
conditionally on a directing variable.
\end{lemma}

\begin{proof}
Suppose a support point had squared norm greater than $\delta$.
Continuity would give a sufficiently small ball about it in which every
pair has inner product greater than $\delta$.  The ball has positive
probability, so independence makes this pair event have positive
probability, contradicting the hypothesis.  A countable base of balls
shows, by separability, that the law is carried by its support.  The
same argument applies to the conditional law for almost every directing
variable.
\end{proof}

Applying \Cref{lem:row-overlap-norm} position by position shows that the
shared vectors of prefix scores satisfy
\begin{equation}
 \max_{h:\,I_h\subset[0,T]}\norm{r_h}\le d^{-b}
 \quad\text{for every fixed }b,
 \label{eq:row-small-shared-score}
\end{equation}
once $d$ is large enough; the exponent in the preceding overlap test
is chosen twice as large.  The other shared vectors need not be small.
Their norms, and all self-Grams, have polynomial bounds in $d$.

We have controlled the shared prefix scores.  To use the private
scores chronologically, we also need the self-Gram identities supplied
by the single-path martingale calculation.  Distinct score intervals
are orthogonal; a score's squared norm equals its bracket in
\Cref{eq:row-bracket-replacement}; and each score is orthogonal to
every centered field observed at or before its interval begins.  For the last assertion, condition on that earlier
field and apply the martingale isometry to its scalar contraction
with the later score.  For score-score and score-square tests, use the
product martingale with the same count stopping as above.  In the
limit these are identities of the self-Gram data, not assertions that
those data have become deterministic.  Moreover
\begin{equation}
 \sum_{h:\,I_h\subset[0,T]} B_h\le \beta^2d,
 \qquad Q_{\mathrm{field},\mathrm{field}}\le4\beta^2.
 \label{eq:row-gram-budgets}
\end{equation}
The first inequality uses $v_j^Y\le1$ and the spatial law of large
numbers for the rate-one counts.  The second uses the deterministic
bound $\norm{Y(s)-\E Y(s)}^2\le4n$.

\subsection{The near-one Gaussian estimate}

The self-Gram identities prepare an adapted Gaussian calculation.
Subtracting the shared Gram leaves the private scores in
\Cref{eq:row-shared-private} with small covariances against earlier
scores and past fields, controlled by \Cref{eq:row-small-shared-score}.
We first correct those covariances to zero.  The correction is made
\emph{conditional on all the self-Gram and bracket data}, which may
still be random.

The correction must also work for a singular covariance matrix.  Set
the forbidden entries to zero, restore the score variances to $B_h$,
and add to the whole matrix a multiple of the identity equal to the
operator norm of that change.  The result is positive semidefinite.
For a fixed $K$, on schedules with at most $Kd$ refreshes, the list has $O_K(d)$
coordinates and its covariance has $O_K(d^2)$ entries. Its perturbation
norm is bounded by the number of coordinates times the largest changed
entry. Thus the added multiple $\delta_d$ can be smaller than any fixed
inverse power of $d$. The final score variances are
$b_h=B_h+\delta_d$, so $\sum_hb_h\le Cd$; replacing the original
compensators $B_h/2$ by $b_h/2$ costs only $O_K(d\delta_d)$.  Couple the original and corrected Gaussians through their
positive square roots.
The inequality
$\norm{Q^{1/2}-\widetilde Q^{1/2}}_{\op}
 \le C\norm{Q-\widetilde Q}_{\op}^{1/2}$
gives Gaussian errors with arbitrarily small polynomial variance.
The common score shifts $g(r_h)$ have the same property.

For each particular spin sequence the change in its coin weight is
bounded multiplicatively by
$\exp(2\sum_h|\Delta H_h|)$, and the change in its likelihood by
$\exp(\sum_h(|\Delta Z_h|+|\Delta B_h|/2))$.
The inequalities survive summation over spin sequences.  Since the
number of positions is $O(d)$, fixed exponential moments of the
comparison errors cost $e^{o(d)}$; increasing the inverse-polynomial
accuracy makes that cost $e^{o(1)}$.  Conditional H\"older inequalities
absorb these errors with an arbitrarily small additional change of
powers.  Thus the correction preserves the desired rate without a
lower bound on the eigenvalues of the covariance matrix.

\begin{lemma}[Chronological Gaussian integration]
\label{lem:row-adapted-gaussian}
Fix self-Gram and bracket data satisfying the corrected identities.
Treat all deterministic mean fields and all common field shifts as
arbitrary forcing.  Let $\widetilde Z_h$ be the corrected prefix
scores, of variances $b_h$, with $\sum_h b_h\le Cd$.
Let $a_h\in\{-1,1\}$ be the spin chosen using the fields and coins
available before score $h$, including a prescribed forced coin when
present.  For
\[
 W=\exp\left\{\sum_h a_h\widetilde Z_h
                         -\tfrac12\sum_h b_h\right\},
\]
one has, for every real $r$,
\begin{equation}
 \begin{aligned}
 \exp\left(-\tfrac12|r(r-1)|Cd\right)
 &\le\E[W^r\mid\text{self-Gram data and forcing}]\\
 &\le\exp\left(\tfrac12|r(r-1)|Cd\right).
 \end{aligned}
 \label{eq:row-near-one}
\end{equation}
In particular, powers within $\eta$ of one cost $e^{O(\eta d)}$.
\end{lemma}

\begin{proof}
Fix the specified data.  The full list is then Gaussian, and the
orthogonality identities make each score independent of the earlier
scores and of every field already used to choose its sign.  Its
variance is $b_h$.  It may remain correlated with fields at later coin
times, which does not affect this chronological calculation.  Generate
the Gaussian list and the coins in time order: the sign $a_h$ is
measurable before revealing $\widetilde Z_h$, so
\[
 \E\left[\exp\{r a_h\widetilde Z_h-r^2b_h/2\}
                      \mid\text{past}\right]=1.
\]
Their product has mean one.  Using $a_h^2=1$ and multiplying by
$\exp\{(r^2-r)\sum_hb_h/2\}$ proves the bounds.  Conditioning on
future self-Gram data is legitimate here because the needed Gaussian
orthogonality identities hold pointwise for each value of those data.
At no point are the scores themselves conditioned on future observed
fields.
\end{proof}

\subsection{Completion of the projected moment estimate}

\begin{proof}[Proof of \Cref{thm:row-projected-moments}]
The prefix is now controlled at powers close to one. Conditional
H\"older inequalities will put every fixed larger power on the short
suffix and on the last-spin factor. First retain schedules with at most
$Kd$ refreshes.  In the limiting representation, factor
\[
 L_\theta F_\sigma=W_{\mathrm{pre}}U,
 \qquad A=\E[W_{\mathrm{pre}}U\mid\mathcal D,g],
\]
where the prefix ends at $T=d-\sqrt d$ and $U$ contains the suffix and
the normalized last-spin factor, wherever that last refresh occurs.
For each fixed $q$, positive and negative moments of the suffix
likelihood are at most $e^{C_q\sqrt d}$.  Apply the forward martingale
calculation on the interval of length $\sqrt d$, and then use the
fixed-$d$ moment passage proved above.  The last-field factor has
bounded moments independent of $d$ in this limiting calculation:
its centered Gaussian field has variance at most $4\beta^2$ by
\Cref{eq:row-gram-budgets}, before the common and private parts are
separated.  H\"older therefore gives
\begin{equation}
 \E[|U|^q+|U|^{-q}]\le 2e^{C_q(1+\sqrt d)}.
 \label{eq:row-suffix-cost}
\end{equation}
The deterministic field means disappear from this bound through
\Cref{eq:row-forced-coin}.

For $0<\eta<1$, conditional H\"older gives the two inequalities
\begin{align}
 A&\le
  \E[W_{\mathrm{pre}}^{1+\eta}\mid\mathcal D,g]^{1/(1+\eta)}
  \E[U^{(1+\eta)/\eta}\mid\mathcal D,g]^{\eta/(1+\eta)},
 \label{eq:row-holder-upper}\\
 \E[W_{\mathrm{pre}}^{1-\eta}\mid\mathcal D,g]
 &\le A^{1-\eta}
       \E[U^{-(1-\eta)/\eta}\mid\mathcal D,g]^\eta.
 \label{eq:row-holder-lower}
\end{align}
For example, the second follows by writing
$W_{\mathrm{pre}}^{1-\eta}
 =(W_{\mathrm{pre}}U)^{1-\eta}U^{-(1-\eta)}$.
This second inequality supplies the inverse-moment bound: inverse
powers of $A$ are controlled by inverse powers of the prefix
conditional average together with positive powers of a suffix moment.

By \Cref{lem:row-adapted-gaussian}, each corrected prefix conditional
average with power $1\pm\eta$ lies between $e^{-C\eta d}$ and
$e^{C\eta d}$, uniformly in the branch data and common field forcing.
Jensen permits integration over the directing law and also taking
inverse powers.  The small common score and covariance-coupling
errors discussed above are handled by H\"older, with a power change
smaller than $\eta$.  Combining
\Cref{eq:row-holder-upper,eq:row-holder-lower,eq:row-suffix-cost} yields
for every fixed $p>0$
\begin{equation}
 \log\E[A^p+A^{-p}]
       \le C_p\eta d+C_{p,\eta}(1+\sqrt d)+o(d).
 \label{eq:row-final-rate}
\end{equation}
This inequality concerns the moments after the dimension limit.  Now
send $d\to\infty$ and then $\eta\downarrow0$.  The rough full-horizon
constants from \Cref{lem:row-uniform-integrability} served only to
identify these limiting moments; they contribute no additional term to
\Cref{eq:row-final-rate}.

H\"older's inequality in \Cref{eq:row-ratio-moment} proves the required
bound for every retained schedule. We next remove the schedule
conditioning. Under the ordinary path law, the tower property gives
\[
 A_{i,x}^M(X_d)
 =\E_x[A_{i,x}^{M,\mathcal S}(X_d)\mid X_d].
\]
The conditional expectation averages the schedule with its ordinary
conditional law given $X_d$. For $r\ge1$ or $r<0$, conditional Jensen
therefore decreases the moment when the schedule is no longer observed.  Control
the omitted schedules by a higher raw likelihood moment, H\"older's
inequality, and the Poisson estimate
\[
 \P\{N_i(d)>Kd\}\le e^{-d(K\log K-K+1)}.
\]
The raw bounds are uniform conditional on the schedule.  Taking $K$
large after the desired moment order makes this contribution smaller
than any prescribed $e^{-Hd}$.  The powers $0\le r\le1$ again follow
from Jensen.  This completes the moment estimate.

It remains to obtain the stated exceptional sets, including for an
integrated family of masks.  Fix $d$, the schedule-count cutoff, the
moment orders, and the comparison accuracies.  Every row concentration
constant above is then uniform in the mask parameters and deterministic
schedule times.  If
$\mathcal B(J,x,\alpha)$ denotes a failure of one of the finitely many
bounded tests, its row-and-start bound has the form
\[
 \E_J\int\mu_J(\dd x)\int
          \1_{\mathcal B(J,x,\alpha)}\,\nu(\dd\alpha)
       \le C e^{-cn}
\]
on the common good cavity events.  Apply Fubini at this point, before choosing the quenched event.  Markov's
inequality gives integrated Gibbs bad mass at most $e^{-cn/2}$ outside
an event of probability $Ce^{-cn/2}$.  If necessary, apply Markov once
more to remove starts with excessive bad-schedule mass, and use higher
rough moments to bound its contribution.  The countable approximations are fixed successively, with
arbitrary fixed slack, so their finite good events suffice before the
dimension limit.  Polynomial endpoint weights and an $e^{o(n)}$
initial density preserve negligibility.  At fixed times the remaining
integrands are dominated by the rough moment bounds and the relevant
mean-square vector budgets.

In a path expansion the order is therefore: fix the global time, hop
cutoff, shapes, elementary caps, moment orders, and label bins; integrate
over their deterministic simplex parameters; take dimension to
infinity; then remove caps and approximation errors; finally increase
the global time and decrease rate slack.  Annulus and first-exceedance
restrictions multiply a positive masked label by factors in $[0,1]$.
Its positive moments are bounded by those of the unrestricted label;
no inverse moment of a restricted label is used.  A selected start or
endpoint vector changes the weights in this integrated assertion and
does not choose an exceptional mask.  Likewise, good interior times
can be chosen from a positive-measure set after an integrated
estimate; a prescribed worst time is not being asserted to be good.
\end{proof}

\begin{corollary}[Stationary endpoint predictors]
\label{cor:row-stationary-predictors}
Let $a_i$ be a nonnegative row label on a span of length $d$, bounded
above by a masked likelihood $\ell_i^M$; factors $v_i\le1$ and cutoff
indicators are allowed.  Its stationary predictor at the right endpoint
has every fixed positive spatial moment bounded by $e^{o(d)}$.
The same bound holds after conditioning at a later endpoint.  Its
predictor at the left endpoint is bounded by one.  Raw fixed-order
moments cost at most $e^{C_pd}$.
\end{corollary}

\begin{proof}
For $p\ge1$, condition on both endpoints of the span and then apply
Jensen to the left endpoint's conditional law given the right endpoint.
The resulting integral over ordinary starts is bounded by
\Cref{thm:row-projected-moments}.  A higher raw moment and H\"older's
inequality make the exponentially small bad-start contribution
negligible.  The Markov property and stationarity make conditioning at
a later endpoint another $L^p$ contraction.  At the left endpoint,
the bound follows directly from
$\E_xa_i\le\E_x\ell_i^M=1$.  Powers between zero and one follow by
concavity.  The raw bound is \Cref{prop:path-rough}.
\end{proof}

\subsection{The positive spin spectral measure}

We turn to the spectral input for hidden propagation.  Let $E_n(\dd u)$
be the spectral resolution of $-L$ on $L^2(\mu)$, and encode the spin
coordinates by the positive matrix-valued measure
\[
 \mathsf M_n(\dd u)_{ij}
      =\ip{x_i}{E_n(\dd u)x_j}_\mu,
 \qquad \rho_n(\dd u)=\frac1n\Tr\mathsf M_n(\dd u).
\]
Spin inversion centers every $x_i$.  Consequently, $\rho_n$ is a
probability measure on the positive spectrum; on the good gap event,
its support lies in $[\gamma,\infty)$.  The heat-bath normalization
also identifies its first two moments:
\begin{equation}
 \begin{aligned}
 \int u\,\rho_n(\dd u)=\frac1n\sum_i\E_\mu v_i\le1,
 \\
 \int u^2\rho_n(\dd u)=\frac1n\sum_i\E_\mu(m_i-x_i)^2
                     =\frac1n\sum_i\E_\mu v_i.
 \end{aligned}
 \label{eq:spin-first-second}
\end{equation}
These bounds give tightness and uniform integrability of the first
moments.  The deterministic stationary spin limits from
\Cref{thm:path-calculus} identify a unique deterministic probability
measure $\rho$ to which the measures converge in probability.  Its
Laplace transform is
\begin{equation}
 C(t)=\int e^{-ut}\rho(\dd u)
     =\lim_{n\to\infty}\frac1n\sum_i\ip{x_i}{S_tx_i}_\mu.
 \label{eq:spin-correlation}
\end{equation}
Set
\[
 \lambda_{\rm sp}=\inf\supp\rho>0,
 \qquad R(t)=-C'(t)=\int u e^{-ut}\rho(\dd u).
\]
For \(z<\lambda_{\rm sp}\), put
\[
 \widehat R(z)=\int_0^\infty e^{zt}R(t)\dd t.
\]
The transform is finite and strictly increasing, and
$\widehat R(0)=1$.  All these definitions take dimension to infinity
first.

\begin{proposition}[Spin edge and the two renewal inequalities]
\label{prop:spin-spectrum}
For $0<\beta<1$,
\begin{equation}
 0<\lambda_{\rm sp}\le(1-\beta)^2,
 \qquad \beta\widehat R(z)<1,
 \qquad \beta^2\widehat R(z)<1-z
 \quad(z<\lambda_{\rm sp}).
 \label{eq:spin-conditions}
\end{equation}
Let $P_0(u)=1$, $P_1(u)=u$, and
$P_{k+1}(u)=uP_k(u)-P_{k-1}(u)$ be the monic second-kind Chebyshev
polynomials.  If
\[
 R_k(t)=\lim_{n\to\infty}\frac1n\Tr\left[
 P_k(J/\beta)\int u e^{-tu}\mathsf M_n(\dd u)\right],
\]
then
\begin{equation}
 \widehat R_k(z)=\beta^k\widehat R(z)^{k+1},
 \qquad z<\lambda_{\rm sp}.
 \label{eq:spin-trunk-transform}
\end{equation}
The limiting covariance trace density with respect to the semicircle
law $\rho_{\rm sc}$ on $[-2,2]$ is
\begin{equation}
 \lim_{n\to\infty}\frac1n
     \Tr\bigl[q(J/\beta)\E_\mu(xx^{\mathsf T})\bigr]
  =\int_{-2}^2\frac{q(u)}{1-\beta u+\beta^2}\rho_{\rm sc}(\dd u)
 \label{eq:spin-covariance-density}
\end{equation}
for every polynomial $q$.
\end{proposition}

\begin{proof}
The proof proceeds from a time-domain diagram identity to the transform
bound, then identifies the covariance trace density and tests its edge.
We first derive the diagram identity, making explicit the role of
stationarity.  Integration by parts and gradient intertwining yield
\begin{equation}
 R(t)=\lim_{n\to\infty}\frac1n\Tr\E_\mu[D_vK_tI].
 \label{eq:spin-gradient-trace}
\end{equation}
Here $K_tI$ applies $K_t$ to each constant coordinate vector.  Insert
the Duhamel expansion and use \Cref{thm:path-calculus}: visible
finite-rank terms vanish in normalized trace, each hop contributes
$\beta sD_v$, and the remaining trace is the sum of complete
noncrossing pairings.  Each resulting vertex carries the
product of its diagonal waiting likelihoods, its arrivals $v$, and
the bare factors $e^{-h}$.  These are nonnegative.  The root initially
carries $v$, as in \Cref{eq:spin-gradient-trace}.

The recurrence for $P_k(s)$ removes exactly the adjacent returning
contraction, making this polynomial the Wick string of $k$ consecutive
edges.  Each of its edges must therefore pair into the Duhamel walk.
Removing those $k$ pairs exposes an ordered trunk of $k+1$ completely
paired rooted time intervals.  Each interval starts with a $v$ weight,
including the interval at the original root, and each trunk edge
contributes $\beta$.  Conversely, concatenate $k+1$ rooted complete
diagrams and adjoin the outside Wick string.  This recovers one and only
one pairing, so no multiplicity factor appears.  After contraction, distinct vertices contribute scalar traces.
Stationarity makes each interval's trace depend only on its duration.
The separation into scalar factors comes from the noncrossing
contraction, rather than from temporal independence of intervals along
one spin history.  The spin/site self-averaging assertion in
\Cref{thm:path-calculus} is what identifies each such scalar factor
with its stationary trace when it is represented by independent
samples.  It is used only for spin, field, and likelihood labels.

The simplex integral over the successive interval lengths therefore
gives the convolution identity
\begin{equation}
 R_k(t)=\beta^k R^{*(k+1)}(t).
 \label{eq:spin-trunk-convolution}
\end{equation}
For the limiting argument, first keep time, Duhamel order, and likelihood
caps fixed.  Then \Cref{prop:path-rough} removes the caps and order
cutoff in the matrix $L^2$ comparison, establishing the identity at
each fixed $t$.  Because the complete-diagram terms are nonnegative,
we may also sum them monotonically when taking Laplace transforms.

To pass to transforms, we need a bound that is valid before the
dimension limit and does not require a finite-dimensional resolvent at
$z$.  Positivity supplies it: for a positive matrix $B$ and a
self-adjoint $Q$, $|\Tr(QB)|\le\norm Q_{\op}\Tr B$.
At finite $n$ apply this to the positive matrix
$\int u e^{-tu}\mathsf M_n(\dd u)$ and to
$Q=P_k(J/\beta)$.  The Wigner spectral bound and
$\sup_{[-2,2]}|P_k|=k+1$ show that, after the dimension limit,
\begin{equation}
 |R_k(t)|\le(k+1)R(t).
 \label{eq:spin-polynomial-domination}
\end{equation}
The same justification can be expressed through truncated integrals:
integrate the finite-dimensional inequality over $0\le t\le T$, take
the dimension limit, and then send $T\to\infty$.  The right-hand side
is integrable for $z<\lambda_{\rm sp}$.  Tonelli's theorem applied to
\Cref{eq:spin-trunk-convolution} now proves
\Cref{eq:spin-trunk-transform} and yields
\[
 (\beta\widehat R(z))^k\le k+1.
\]
Taking $k$th roots yields $\beta\widehat R(z)\le1$.  Given
$z<\lambda_{\rm sp}$, choose $z'$ strictly between them.  Since
$\widehat R$ is strictly increasing,
$\beta\widehat R(z)<\beta\widehat R(z')\le1$.

Integrating \Cref{eq:spin-trunk-convolution} at $z=0$ gives
\[
 \lim_{n\to\infty}\frac1n\Tr\left[
   P_k(J/\beta)\E_\mu(xx^{\mathsf T})\right]=\beta^k.
\]
Once again, justify the integral by first stopping at fixed $T$.
The spectral gap bounds the omitted covariance tail by
$\norm{P_k(J/\beta)}_{\op}e^{-\gamma T}$.  To identify the resulting
covariance density, use the orthonormality of the $P_k$ for
$\rho_{\rm sc}$ and the generating function
\[
 \sum_{k\ge0}\beta^kP_k(u)=\frac1{1-\beta u+\beta^2}
 \quad(-2\le u\le2);
\]
the series converges uniformly because $\beta<1$ and
$|P_k(u)|\le k+1$.  This proves
\Cref{eq:spin-covariance-density}.  The measure identified here is
positive: at finite $n$ it is the trace of the positive covariance
matrix against the spectral projections of $J/\beta$.

We next use this covariance density to bound the spin edge.  Fix a real
polynomial $q$.  Weighting $\mathsf M_n$ by $q(J/\beta)^2$ gives a
nonnegative scalar spectral measure, dominated by
$\norm{q(J/\beta)}_{\op}^2\rho_n$.  All its limiting mass therefore
lies in $[\lambda_{\rm sp},\infty)$, even though the finite-dimensional
spin measures may put small masses below that edge.  Its first moment
satisfies
\[
 \frac1n\Tr\bigl[q(J/\beta)^2D_{\E_\mu v}\bigr]
       \le\frac1n\Tr q(J/\beta)^2,
\]
because the energy matrix
$\int u\mathsf M_n(\dd u)$ is exactly $D_{\E_\mu v}$.
Use \Cref{eq:spin-covariance-density}, and then the semicircle trace
limit, to obtain
\begin{equation}
 \lambda_{\rm sp}
   \int\frac{q(u)^2}{1-\beta u+\beta^2}\rho_{\rm sc}(\dd u)
       \le\int q(u)^2\rho_{\rm sc}(\dd u).
 \label{eq:spin-edge-test}
\end{equation}
Alternatively, discard scalar spectral mass below
$\lambda_{\rm sp}-\delta$ before sending $\delta\downarrow0$.  This
gives the same inequality without requiring a uniform finite-$n$ lower
edge.
Choose $q_m(u)=((u+2)/4)^m$.  After normalization, the measures
$q_m^2\rho_{\rm sc}$ concentrate at $2$.  Letting $m\to\infty$ in
\Cref{eq:spin-edge-test} gives
$\lambda_{\rm sp}\le1-2\beta+\beta^2$.

Finally, for $0\le z<\lambda_{\rm sp}$,
\[
 \beta^2\widehat R(z)<\beta
       <2\beta-\beta^2
       \le1-\lambda_{\rm sp}<1-z.
\]
For $z<0$, $\widehat R(z)\le\widehat R(0)=1$ and
$\beta^2<1<1-z$.  These prove the second renewal inequality and
complete \Cref{eq:spin-conditions}.
\end{proof}

The two strict inequalities control different vertices in the path
expansion.  An ordinary interior trunk vertex has integrated weight
$\beta\widehat R(z)<1$.  If its initial variance factor is dropped,
we instead retain the bare waits and sum any number of paired side
excursions.  The resulting bound is the geometric renewal resolvent
\begin{equation}
 \frac1{1-z-\beta^2\widehat R(z)}<\infty.
 \label{eq:spin-exceptional-renewal}
\end{equation}
A bare wait contributes the transform $(1-z)^{-1}$.  Each inserted
excursion contributes two hops and one rooted paired interval, hence
the transform $\beta^2\widehat R(z)$.  Summing successive insertions
gives \Cref{eq:spin-exceptional-renewal}, with the bare waits retained
throughout.  The bound applies because $z<\lambda_{\rm sp}<1$.
